\chapter{Helium Restricted Hartree--Fock Control}

\section{Self-consistent field equation}
The one-parameter $1s^2$ ansatz establishes the first interacting variational bound but does not permit the orbital to relax freely. The next control layer solves the spherical restricted Hartree--Fock equation numerically. For a helium-like two-electron ion with opposite spins sharing one spatial orbital $\phi$, the Fock equation is
\begin{equation}
\left[-\frac12\nabla^2-\frac{Z}{r}+J[\phi](r)\right]\phi(r)=\varepsilon\phi(r),
\end{equation}
where the self-interaction is cancelled by exchange and the remaining Coulomb field is that of the opposite-spin electron.

Writing the $s$ orbital as $\phi(r)=u(r)/(r\sqrt{4\pi})$ with $\int_0^\infty |u(r)|^2\,dr=1$, spherical symmetry reduces the Coulomb potential to
\begin{equation}
J(r)=\frac{1}{r}\int_0^r |u(r')|^2\,dr'+\int_r^\infty\frac{|u(r')|^2}{r'}\,dr'.
\end{equation}
The total RHF energy is evaluated as
\begin{equation}
E_{\rm RHF}=2\langle\phi|h|\phi\rangle+J,
\end{equation}
with $h=-\nabla^2/2-Z/r$.

\section{Discretization and SCF}
The implementation uses a uniform radial grid, a second-order three-point finite-difference kinetic operator, direct spherical Coulomb quadrature, density/orbital mixing, and repeated diagonalization of the one-orbital Fock operator until the total energy changes by less than the requested tolerance. No experimental energy is used inside the SCF loop.

Because the finite-difference error is second order in the grid spacing, two nested grids with steps $h$ and $h/2$ are combined by Richardson extrapolation,
\begin{equation}
E_{\infty}\approx\frac{4E_{h/2}-E_h}{3}.
\end{equation}

\section{Helium benchmark}
For $Z=2$, $r_{\max}=20a_0$, and nested grids of 399 and 799 interior points, the implementation obtains
\begin{align}
E_{399} &= -2.8534029917\ {\rm Ha},\\
E_{799} &= -2.8596016563\ {\rm Ha},\\
E_{\rm Richardson} &= -2.8616678778\ {\rm Ha}.
\end{align}
Against the diagnostic Hartree--Fock limit $-2.861679995612$ Ha, the relative discrepancy is approximately $4.23\times10^{-6}$, or $0.000423\%$. Both SCF calculations converge in 19 iterations under the default high-accuracy gate.

\section{Interpretation and TIR boundary}
This result closes the first mean-field control requirement for the atom programme: the code can recover the conventional helium RHF limit without any TIR correction. Therefore the missing correlation energy must not be relabelled as ``resonance'' or as an informational interaction. The next control layer is an explicitly correlated or configuration-interaction calculation. Only after conventional correlation is quantified will a TIR relational observable be introduced, and it must predict an independent invariant or residual rather than repair a known energy by fitting.
